% Source: Weinberg GR, physical PDF pp. 202--205
%         (printed pp. 179--182).
% Coverage: Section 8.2, The Schwarzschild Solution; equations
%           (8.2.1)--(8.2.16).
% Figures/tables/footnotes: none.
% Status: source-reviewed and compile-clean.
% Uncertainties: Ricci signs in (8.2.3), (8.2.7), and (8.2.8) converted
%                consistently with the binding curvature convention.

\subsection{The Schwarzschild Solution}\label{sec:8.2}

We now apply the Einstein field equations to the general static isotropic
metric. We use the standard form discussed in the last section, that is,
\begin{equation}
  ds^2
  =
  -B(r)\,dt^2+A(r)\,dr^2
  +r^2d\theta^2+r^2\sin^2\theta\,d\varphi^2.
  \tag{8.2.1}\label{eq:8.2.1}
\end{equation}
The field equations for empty space are
\begin{equation}
  R_{\mu\nu}=0.
  \tag{8.2.2}\label{eq:8.2.2}
\end{equation}
The components of the Ricci tensor are given for this metric by
Eq.~\eqref{eq:8.1.13}. We see that it will suffice to set
\(R_{rr}\), \(R_{\theta\theta}\), and \(R_{tt}\) equal to zero. We also
see that
\begin{equation}
  \frac{R_{rr}}{A}+\frac{R_{tt}}{B}
  =
  \frac{1}{rA}\left(\frac{A'}{A}+\frac{B'}{B}\right),
  \tag{8.2.3}\label{eq:8.2.3}
\end{equation}
% MODERNIZATION CHECK: the right-hand side is the negative of Weinberg's
% printed (8.2.3), because both Ricci components use the target convention.
so Eq.~\eqref{eq:8.2.2} requires that \(B'/B=-A'/A\), or
\begin{equation}
  A(r)B(r)=\text{constant}.
  \tag{8.2.4}\label{eq:8.2.4}
\end{equation}
Furthermore, we impose on \(A\) and \(B\) the boundary condition that for
\(r\to\infty\) the metric tensor must approach the Minkowski tensor in
spherical coordinates, that is,
\begin{equation}
  \lim_{r\to\infty}A(r)=\lim_{r\to\infty}B(r)=1.
  \tag{8.2.5}\label{eq:8.2.5}
\end{equation}
From Eqs.~\eqref{eq:8.2.4} and \eqref{eq:8.2.5} we have then
\begin{equation}
  A(r)=\frac{1}{B(r)}.
  \tag{8.2.6}\label{eq:8.2.6}
\end{equation}
Since Eq.~\eqref{eq:8.2.3} now vanishes, it remains to make
\(R_{tt}\) and \(R_{\theta\theta}\) vanish. Using Eq.~\eqref{eq:8.2.6}
in Eq.~\eqref{eq:8.1.13}, we find
\begin{equation}
  R_{\theta\theta}=1-rB'(r)-B(r),
  \tag{8.2.7}\label{eq:8.2.7}
\end{equation}
\begin{equation}
  R_{tt}
  =
  \frac{B(r)B''(r)}{2}
  +\frac{B(r)B'(r)}{r}
  =
  \frac{B(r)}{2r}\bigl(rB(r)\bigr)'',
  \tag{8.2.8}\label{eq:8.2.8}
\end{equation}
so it is sufficient to set \(R_{\theta\theta}\) equal to zero, that is,
\[
  \frac{d}{dr}\bigl(rB(r)\bigr)=rB'(r)+B(r)=1.
\]
The solution is
\begin{equation}
  rB(r)=r+\text{constant}.
  \tag{8.2.9}\label{eq:8.2.9}
\end{equation}
To fix the constant of integration we recall that at great distances from a
central mass \(M\), the component \(g_{tt}=-B\) must approach
\(-1-2\phi\), where \(\phi\) is the Newtonian potential \(-MG/r\).
(See Section~\ref{sec:3.4}.) Hence the constant of integration is
\(-2MG\), and our final solution is
\begin{equation}
  B(r)=1-\frac{2MG}{r},
  \tag{8.2.10}\label{eq:8.2.10}
\end{equation}
\begin{equation}
  A(r)=\left(1-\frac{2MG}{r}\right)^{-1}.
  \tag{8.2.11}\label{eq:8.2.11}
\end{equation}
The full metric is given by
\begin{equation}
  ds^2
  =
  -\left(1-\frac{2MG}{r}\right)dt^2
  +\left(1-\frac{2MG}{r}\right)^{-1}dr^2
  +r^2d\theta^2+r^2\sin^2\theta\,d\varphi^2.
  \tag{8.2.12}\label{eq:8.2.12}
\end{equation}
This solution was found by K.~Schwarzschild in 1916.

The Schwarzschild solution is expressed in Eq.~\eqref{eq:8.2.12} in its
``standard'' form. We can also express it in the equivalent ``isotropic''
form, by introducing a new radius variable
\begin{equation}
  \rho
  \equiv
  \frac12\left[r-MG+\left(r^2-2MGr\right)^{1/2}\right],
  \tag{8.2.13}\label{eq:8.2.13}
\end{equation}
or
\[
  r=\rho\left(1+\frac{MG}{2\rho}\right)^2.
\]
Substituting this in Eq.~\eqref{eq:8.2.12} gives
\begin{equation}
  \begin{split}
  ds^2={}&
  -\frac{(1-MG/2\rho)^2}{(1+MG/2\rho)^2}\,dt^2\\
  &+\left(1+\frac{MG}{2\rho}\right)^4
  \left(d\rho^2+\rho^2d\theta^2
  +\rho^2\sin^2\theta\,d\varphi^2\right).
  \end{split}
  \tag{8.2.14}\label{eq:8.2.14}
\end{equation}
We can also construct harmonic coordinates
\[
  X_1=R\sin\theta\cos\varphi,\qquad
  X_2=R\sin\theta\sin\varphi,\qquad
  X_3=R\cos\theta,\qquad t,
\]
by using for \(R\) a solution of the differential equation
\eqref{eq:8.1.15}, which here becomes
\[
  \frac{d}{dr}
  \left[
    r^2\left(1-\frac{2MG}{r}\right)\frac{dR}{dr}
  \right]-2R=0.
\]
One convenient solution is
\[
  R=r-MG.
\]
The metric is then given by Eq.~\eqref{eq:8.1.16}:
\begin{equation}
  \begin{split}
  ds^2={}&
  -\left(\frac{1-MG/R}{1+MG/R}\right)dt^2
  +\left(1+\frac{MG}{R}\right)^2d\mathbf{X}^2\\
  &+\left(\frac{1+MG/R}{1-MG/R}\right)
  \frac{M^2G^2}{R^4}
  (\mathbf{X}\cdot d\mathbf{X})^2,
  \qquad R^2\equiv\mathbf{X}^2.
  \end{split}
  \tag{8.2.15}\label{eq:8.2.15}
\end{equation}

We identified the integration constant \(M\) with the mass of the sun by
comparison with Newton's theory. In fact, we can show that \(M\) is
precisely equal to the total energy \(P^0\) of the sun \emph{and its
gravitational field}. Let us write the standard form of the metric in
quasi-Minkowskian coordinates, by defining
\[
  x^1=r\sin\theta\cos\varphi,\qquad
  x^2=r\sin\theta\sin\varphi,\qquad
  x^3=r\cos\theta.
\]
Then Eq.~\eqref{eq:8.2.12} becomes
\[
  ds^2
  =
  -\left(1-\frac{2MG}{r}\right)dt^2
  +\left\{
    \left(1-\frac{2MG}{r}\right)^{-1}-1
  \right\}r^{-2}(\mathbf{x}\cdot d\mathbf{x})^2
  +d\mathbf{x}^2.
\]
Since \(g_{\mu\nu}\) is time independent and \(g_{i0}\) vanishes, it
follows from Eq.~\eqref{eq:7.6.22} that the total momentum \(P^i\) of the
system vanishes, which of course it must do since the system is static and
isotropic. To calculate the total energy, we need the asymptotic behavior of
the spatial part of the metric, as \(r\to\infty\),
\[
  h_{ij}\equiv g_{ij}-\delta_{ij}
  \longrightarrow
  \frac{2MG}{r}n_i n_j+\order(r^{-2}),
\]
where \(n_i=x^i/r\). To calculate the integral \eqref{eq:7.6.21} we use
the relations
\[
  \partial_i r=n_i,
  \qquad
  \partial_j n_i=\frac{\delta_{ij}-n_i n_j}{r},
\]
and find
\[
  \partial_i h_{jj}-\partial_j h_{ij}
  \longrightarrow
  -\frac{4MG}{r^2}n_i+\order(r^{-3}),
\]
so Eq.~\eqref{eq:7.6.23} gives the total energy of matter and gravitation
here as
\begin{equation}
  P^0=M.
  \tag{8.2.16}\label{eq:8.2.16}
\end{equation}
The reader may check that the same result would be given by the isotropic or
harmonic forms of the Schwarzschild solution. Finally,
Eq.~\eqref{eq:7.6.24} gives for the total angular momentum here the
expected value zero.
